\section{Governance and Organisation}



\begin{frame}[allowframebreaks]{Project Kickoff}
    \mytext{
    First kickoff... then kickoff.. then... wait?
    \medskip
    What exactly is a `kickoff'?
    \medskip
    \begin{itemize}
        \item Rewarding of contract
        \item Celebrating closure of contract 
        \item Public declaration
        \item Ground breaking ceremony
        \item Topping out ceremony
    \end{itemize}
    }
\end{frame}
    
    
    
    
%     \begin{itemize}
%         \item define clear milestones
%         \item define clear outcomes of meetings
%         \item make decisions public
%     \end{itemize}
% \end{frame}


% \begin{frame}[allowframebreaks]{Coming up with a Project}  % depends on: smartdiagram, adjustbox
%     \begin{adjustbox}{width=\textwidth}
%         \smartdiagramset{back arrow disabled=true}
%         \smartdiagram[flow diagram:horizontal]{Notice, Understand, Find, Formulate, Convince}
%     \end{adjustbox}
%     \begin{description}
%         \item[Notice]  {a problem / an opportunity},
%         \item[Understand] {causes / the situation}, 
%         \item[Find] {a solution / a business case}, 
%         \item[Formulate] {the business / project plan},
%         \item[Convince] {stakeholders},
%         \item[Plan] {tasks, budget, team, time},
%         \item[Do] {work, controlling, reporting, correcting}
%         \item[Deliver] {product, result, review}
%     \end{description}
    
%     \framebreak
    
%     \begin{adjustbox}{width=\textwidth}
%         \smartdiagramset{back arrow disabled=true}
%         \smartdiagram[flow diagram:horizontal]{receive notice, analyse, formulate, begin}
%     \end{adjustbox}
    

%     \begin{description}
%         \item[Receive notice]  {from project sponsor},
%         \item[Understand] {stakeholders, their role, motivation etc.}, 
%         \item[Find] {a strategy}, 
%         \item[Formulate] {the project plan},
%         \item[Convince] {stakeholders},
%         \item[Do] {planning, work, controlling, reporting, correcting}
%         \item[Deliver] {product, result, review}
%     \end{description}

% \end{frame}




% \begin{frame}[allowframebreaks]{Typical Main Project Phases (Suchanek)}
%     \resizebox{!}{0.8\textheight}{
%         \input{tikz/tikz-chain-projectphases-details}
%     }

%     \framebreak
    
%     \resizebox{\textwidth}{!}{
%     \begin{meinganttq}{5}
%         \inserttable{gantts/mains.csv}
%         \showmilestones{gantts/mains.csv}
%     \end{meinganttq}
%     }
% \end{frame}

\begin{frame}{Project Launch per Kerzner}


\mytext{
    In the launch phase, customer (= project sponsor), project manager (pm), subject matter experts (sme's) and functional leads (leads) are involved. 
    \begin{description}
    \item[Phase 1] - Initialisation
    \begin{itemize}
        \item Kickoff-Meeting with customer
    \end{itemize}
    \item[Phase 2] - Preliminary planning
    \begin{itemize}
        \item PM, SMEs, Leads: Scope Management plan
        \item Preliminary Scope Statement
    \end{itemize}
    \end{description}
    
}
\end{frame}

\begin{frame}{Project Launch per Kerzner}

\mytext{
    \begin{description}
    \item[Phase 3] - Detailed planning
    \begin{itemize}
        \item PM, SMEs, Leads: \textbf{Work Breakdown Structure (WBS)} development
        \item WBS dictionary
    \end{itemize}
    \item[Phase 4] Scope baseline
    \begin{itemize}
        \item Finalisation of governing documents\\ ("Lastenheft", product requirement documentation)
    \end{itemize}
    \end{description}
    
    Source: \cite[p. 158]{KerznerHarold2017Pm:a}
}
\end{frame}

% \begin{frame}{Project Kick off}
% Example: project `play football'
%     \begin{itemize}
%         \item Before you kick the ball...  
%         \item Players must already be committed
%         \item Rules and commitment must be clear
%         \item You don't want an embarrassment
%         \item Your ideas must already be their ideas
%         \item Stakeholders are a team
%         \item Ceremonial event, team building
%         \item Believe in success!
%     \end{itemize}
% \end{frame}



% \begin{frame}{The Objectives of a Project}

%     Well defined objectives are prerequisite. BUT
%     \begin{itemize}
%         \item objectives may be subject to change.
%         \item different stake holders may have different views.
%     \end{itemize}
    
    
%     Defined roles and institutions are much less easy to change.
    
% \end{frame}


% \begin{frame}{Project Governance Document (PGD)}
%     \begin{itemize}
%         \item A PGD is also the `constitution' of the project
%         \item you must identify / find stakeholders
%         \item you have to create institutions
%         \item define obligations / duties
%     \end{itemize}
% \end{frame}

% \begin{frame}{From Knowledge to Action}

%     \begin{quote}
%         \flushright Der Mensch kann viel, wenn er nur muss\\
%         \flushright---unknown
%     \end{quote}

%     \begin{itemize}
%         \item to have the right to
%         \item to know how to
%         \item to have reason to
%         \item to want to
%         \item to have what it takes to
%         \item to have the duty to
%     \end{itemize}
% \end{frame}

% \begin{frame}{Obligations / Duties}
%     \begin{itemize}
%         \item bind yourself to your word
%         \item bind yourself to your supporters
%         \item bind yourself to your goals
%         \item set yourself s.m.a.r.t. goals
%         \item promise stakeholders s.m.a.r.t. goals
%         \item set team members and contractors s.m.a.r.t. goals
%         \item report and control
%     \end{itemize}
% \end{frame}


% \begin{frame}{Idea, initialisation}
%     \begin{itemize}
%         \item Why do it?
%         \item What to do to achieve what?
%         \item Who wants it?
%         \item Who pays for it?
%         \item Who might oppose it?
%         \item What will we need?
%         \item (What are the opportunities?)
%         \item (What are the risks?)
%         \item (What are the drawbacks?)
%         \item Whom do we need to convince?
%         \item To whom do we have to communicate?
%     \end{itemize}
% \end{frame}

% \begin{frame}{What Counts}
%     \begin{itemize}
%         \item Relevant for the Project Owner / Sponsor:
%         \begin{itemize}
%             \item opportunities and risks of the project
%             \item business case, business strategy
%             \item \textbf{why}
%         \end{itemize}
%         \item For the Project Manager:
%         \begin{itemize}
%             \item \textbf{who} wants this? 
%             \item \textbf{who} does it?
%             \item \textbf{what do they actually want?}
%             \item \textbf{with what }can we do it?
%             \item opportunities and risks \textbf{of tasks / sub-projects}
%             \item stakeholder satisfaction
%             \item formulate s.m.a.r.t.~goals
%             \item \textbf{how} we can do it
%             \item \textbf{when} will we do it
%         \end{itemize}
%     \end{itemize}
% \end{frame}



\begin{frame}{The Double Perspective of the Project Manager}

    \showmytikz{../shared/tikz/tikz-doubleperspective}

\end{frame}



% \begin{frame}[allowframebreaks]{Stake holders and Management}
  
%       \textbf{Someone} wants \textbf{something} to happen 
    
%     \begin{itemize}
%         \item Understand who wants the project to happen
%         \item What does this someone give you?
%         \item What does this someone expect?
%         \item What is this something?
%         \item How can it be achieved?
%         \item What do we know, and what do we not know?
%         \item How do we communicate and, if required, correct course?
%         \item Who may intervene? How do you react?
%     \end{itemize}
  
%     \framebreak 
    
%     Complex projects have \textbf{many stake holders.} 
    
%     \begin{itemize}
%         \item inside vs. outside
%         \item powerful vs. not powerful
%     \end{itemize}
    
    
%     $\rightarrow$ `Stakeholder Matrix'
%     \framebreak
    
%     \begin{itemize}
%         \item `above', such as project sponsors, project steering team, top level management and authorities, customers...
%         \item `below', such as the managers of sub-teams and  executing bodies
%         \item `on the same level', such as supporting line functions within the company
%         \item `outside' such as press etc.
%     \end{itemize}


%     \begin{tikzpicture}
%         \draw[<->] (-1.5,0) -- (1.5,0);
%         \draw[<->] (0,-1.5) -- (0,1.5);
%         \node[circle,draw,fill=white] at (0,0) {PM};
%     \end{tikzpicture}

%     Communication both to identify plans and to inform stake holders is a large part of project management.
% \end{frame}



\begin{frame}{Organisation and Governance}
    \mytext{ 
        \setlength{\parskip}{3ex}
        Governance is the system by which an organisation is controlled and operates. 
        
        Projects are one-time efforts: it follows that the organisation is created for the occasion. 
        
        But how are stakeholders, managers, partners etc. organised?
        
        
    }
    
\end{frame}    

\begin{frame}{Organisation and Governance}
    \mytext{ 
        \setlength{\parskip}{3ex}    
        Typically, a company is founded just for the purpose of the project (Special Purpose Vehicle or \textbf{SPV}).
        
        Various legal forms are possible. 
        
        Setting up, managing and evolving the organisation and its rules is a core task of project management. 
    
    \begin{itemize}
        \item In-house project: setting up a team 
        \item Otherwise: a company / association
    \end{itemize}
    
    }
    
\end{frame}    



\begin{frame}{From Project to Product: `Port'}
    \begin{hsbfull}
        \showmytikz[caption={From Project to Product: `Port' (Stakeholders)},label=fig:port-stakeholders]{../shared/tikz/tikz-port.tex}
    \end{hsbfull}
\end{frame}



\begin{frame}{Organisation and Governance}
    \mytext{
    \begin{itemize}
        \item Organise Relationships
        \item Define roles and responsibilities
        \item Decision making structure
        \item Steering the project
        \item Escalation management
        \item Change management
        \item Communication Strategy
    \end{itemize}
    }
\end{frame}

\begin{frame}{Projects within a Company}
    \showmytikz[1]{../shared/tikz/tikz-projectswithinacompany}
\end{frame}

\begin{frame}{Projects within a Company}
    \showmytikz[1,2]{../shared/tikz/tikz-projectswithinacompany}
\end{frame}

\begin{frame}{Projects within a Company}
    \showmytikz[2,3]{../shared/tikz/tikz-projectswithinacompany}
\end{frame}



\begin{frame}{Project Oriented Organisation}
    \showmytikz[1]{../shared/tikz/tikz-project-oriented-organisation}
\end{frame}

\begin{frame}{Project Oriented Organisation}
    \showmytikz[2]{../shared/tikz/tikz-project-oriented-organisation}
\end{frame}

\begin{frame}{Project Oriented Organisation}
    \showmytikz[1,2,3,4,5]{../shared/tikz/tikz-project-oriented-organisation}
\end{frame}




\begin{frame}{Hierarchy}
    \showmytikz{../shared/tikz/tikz-PM-smartd-hierarchy}
    \mycomment{
    

    \vfill
    The reality is more complicated!
    }
\end{frame}

\begin{frame}{Core Team Functions}
    \mytext{
        \begin{itemize}
            \item Lead(er), Manager (Product Owner?)
            \item Assistant
            \item (financial) Administration
            \item IT (support)
            \item Legal
            \item Subject Matter Experts 
            \item Partners
            \item ...
            \item Cousellor (Scrum master?)
        \end{itemize}
        
        Different ways to organise team work,
        
        different philosophies...
        
        (to be continued)
    }
\end{frame}

\begin{frame}{Governance / Rules}
    \mytext{
    
    \begin{itemize}
        \item every team should have rules
        \item typically not democratic
        \item consensus, decisions
        \item degree of rules may vary
        \item structure and clarity
        \item (psychological) safety
    \end{itemize}
    
    \medskip
    (to be continued)
    
    \medskip
    But: the team itself is also not alone!
    }
    
\end{frame}


\begin{frame}{Constitution}
    \showmytikz{../shared/tikz/tikz-constitution.tex}
\end{frame}

\begin{frame}{Project Constitution (Organisation Chart)}
    \showmytikz[caption={Project Constitution (Governance)},label=fig:projectconstitution-governance]{../shared/tikz/tikz-projectconstitution}
\end{frame}

% \begin{frame}{From Project to Product: `Port'}
%     \showmytikz[caption={From Project to Product: `Port' (Governance)},label=fig:port-governance]{../shared/tikz/tikz-port.tex}
% \end{frame}


% \begin{frame}{Project Governance Document (PGD)}
%     \begin{itemize}
%         \item A PGD is also the `constitution' of the project
%         \item you must identify / find stakeholders
%         \item you have to create institutions
%         \item define obligations / duties
%     \end{itemize}
% \end{frame}

% \begin{frame}{From Knowledge to Action}

%     \begin{quote}
%         \flushright Der Mensch kann viel, wenn er nur muss\\
%         \flushright---unknown
%     \end{quote}

%     \begin{itemize}
%         \item to have the right to
%         \item to know how to
%         \item to have reason to
%         \item to want to
%         \item to have what it takes to
%         \item to have the duty to
%     \end{itemize}
% \end{frame}

% \begin{frame}{Obligations / Duties}
%     \begin{itemize}
%         \item bind yourself to your word
%         \item bind yourself to your supporters
%         \item bind yourself to your goals
%         \item set yourself s.m.a.r.t. goals
%         \item promise stakeholders s.m.a.r.t. goals
%         \item set team members and contractors s.m.a.r.t. goals
%         \item report and control
%     \end{itemize}
% \end{frame}


% \begin{frame}{Idea, initialisation}
%     \begin{itemize}
%         \item Why do it?
%         \item What to do to achieve what?
%         \item Who wants it?
%         \item Who pays for it?
%         \item Who might oppose it?
%         \item What will we need?
%         \item (What are the opportunities?)
%         \item (What are the risks?)
%         \item (What are the drawbacks?)
%         \item Whom do we need to convince?
%         \item To whom do we have to communicate?
%     \end{itemize}
% \end{frame}

% \begin{frame}{What Counts}
%     \begin{itemize}
%         \item Relevant for the Project Owner / Sponsor:
%         \begin{itemize}
%             \item opportunities and risks of the project
%             \item business case, business strategy
%             \item \textbf{why}
%         \end{itemize}
%         \item For the Project Manager:
%         \begin{itemize}
%             \item opportunities and risks \textbf{of tasks / sub-projects}
%             \item stakeholder satisfaction
%             \item formulate s.m.a.r.t.~goals
%             \item how
%             \item when
%             \item \textbf{who}
%             \item \textbf{with what}
%         \end{itemize}
%     \end{itemize}
% \end{frame}



% \begin{frame}{The Double Perspective of the Project Manager}
%     \begin{tikzpicture}
%         \draw  (0,0) node[below right] {Project Owner, Project Sponsors, Stake Holders} rectangle (10,-1);
        
%         \draw[->] (1,-1) -- node[right] {is entitled to $P$ and wants $P$} (1,-3);
%         \draw[<-] (6,-1) -- node[right] {is obliged to deliver $P$} (6,-3);
        
%         \draw  (0,-3) node[below right] {Project Manager w/ Steering Group} rectangle (10,-4);
        
%         \draw[->] (1,-4) -- node[right] {is entitled to $p$ and wants $p$} (1,-6);
%         \draw[<-] (6,-4) -- node[right] {is obliged to deliver $p$} (6,-6);
        
%         \draw  (0,-6) node[below right] {Project Management Office (PMO), contractors} rectangle (10,-7);
%     \end{tikzpicture}
% \end{frame}

\begin{frame}{An Example: Public Sector Project}
    \mytext{
    Public Sector Projects are typically financed by the government as part of some larger programme. The ministry in charge then contracts a company (itself possibly in public ownership, but not necessarily so) to administer the programme's fund, to reward successful project applications, to supervise the projects, and to administer.

    Locally, consortiums send in proposals for projects which then get rewarded contract. The result is a multi-layered governance structure.
    }   
\end{frame}

\begin{frame}{An Example: Public Sector Project}
    \showmytikz{../shared/tikz/tikz-publicsectorproject}
\end{frame}

\begin{frame}{Example: Establishing an Organisation}
    \mytext{
    Not all projects produce tangible results, but all are governed by some organisation. Setting up an organisation is thus part of \textit{every }project. 

    In our next example, the aim of the project is to establish an association of some members.
    }
\end{frame}


\begin{frame}{Example: Establishing an Organisation}
    \showmytikz{../shared/tikz/tikz-association}
\end{frame}


\begin{frame}{Initiators, Members, Partners...}
    \mycomment{
    In such a setup, where multiple legal and natural persons are involved, a clear vocabulary should be established. Everyone is a partner. But not everyone is in the consortium, and not everyone a member.

    Not only does this help to make clear whom we are talking about in a given context, but we can also decide on rules for each of these groups.    
    }
    \showmytikz{../shared/tikz/tikz-initiatorsmemberspartners}
\end{frame}

\begin{frame}{PMO and Steering Group}
    \mytext{
        The steering group should mirror in composition the teams involved in the project: as below, so above.

        This could be by personal identity, or could be by subject matter experts.
    }   
\end{frame}

\begin{frame}{PMO and Steering Group}
    \begin{hsbfull}
        \showmytikz[1]{../shared/tikz/tikz-pmoandsteeringgroup}
    \end{hsbfull}
\end{frame}

\beameronly{
    \begin{frame}{PMO and Steering Group}
        \begin{hsbfull}
            \showmytikz[1,2]{../shared/tikz/tikz-pmoandsteeringgroup}
        \end{hsbfull}
    \end{frame}
    
    \begin{frame}{PMO and Steering Group}
        \begin{hsbfull}
            \showmytikz[1,2,3]{../shared/tikz/tikz-pmoandsteeringgroup}
        \end{hsbfull}
    \end{frame}
}



\begin{frame}{Governance and Work Packages}
    \mycomment{
        In complex projects, the steering board should be composed of representatives of those initiators who also take over work packages of the project. Governance and project planning / work breakdown is thus intrinsically connected.

        This way, everyone involved in the  \enquote{doing} is also represented on the policy setting level, and the policy setting level knows about the challenges of ground work done.
    }
        \showmytikz{../shared/tikz/tikz-projectgovernance}
\end{frame}


% \begin{frame}{Project Phases}
%     \resizebox{!}{0.8\textheight}{
%         \input{tikz/tikz-chain-projectphases-details}
%     }
% \end{frame}
